\documentclass[12pt]{article}
\usepackage[T1]{fontenc}
\usepackage[utf8]{inputenc}
\usepackage{lmodern}
\usepackage{textcomp}
\usepackage{enumitem}
\usepackage{geometry}
\geometry{margin=1in}
\begin{document}
\begin{center}
Answer Key - Philosophy - Undergraduate Level Assessment 4 \
\small (Answers are ordered exactly as in the question paper.)
\end{center}

\section*{Multiple Choice}
\begin{enumerate}[label=\arabic*.]
\item \textbf{Answer:} E
\\ \textit{Options:} A) disagreement in principles.; B) disagreement in understanding.; C) disagreement in experience.; D) disagreement in practice.; E) disagreement in interest.
\\ \textit{Note:} According to Stevenson, moral disagreements arise from differing interests and values, which highlights that such disagreements are fundamentally rooted in conflicting desires rather than objective truths.
\item \textbf{Answer:} B
\\ \textit{Options:} A) seems to be right, although this might be illusory.; B) tends to be right, although this might be outweighed by other considerations.; C) is right.; D) is the first thing that an agent ought to do, above all else.
\\ \textit{Note:} A prima facie duty refers to the moral obligation that an action appears to have in favor of its rightness, but this obligation can be overridden by more pressing moral considerations in specific situations.
\item \textbf{Answer:} A
\\ \textit{Options:} A) educational institutions.; B) environmental entities.; C) animals.; D) governments.; E) corporations.
\\ \textit{Note:} Sinnott-Armstrong focuses on the moral responsibilities of various entities, and educational institutions play a crucial role in shaping ethical understanding and societal values, making them a primary area of his inquiry.
\item \textbf{Answer:} A
\\ \textit{Options:} A) cause behavior; B) are the same as behavioral states; C) exist; D) can be directly observed; E) are always visible
\\ \textit{Note:} Logical behaviorism posits that mental states can be fully understood in terms of observable behaviors, which contradicts the commonsense intuition that mental states independently cause behaviors.
\item \textbf{Answer:} E
\\ \textit{Options:} A) Mism; B) Mssi; C) Msim; D) Mmis; E) Mmsi
\\ \textit{Note:} The correct answer Mmsi accurately represents the action of Marco (m) moving from Italy (i) to Spain (s) by using the predicate Mxyz, which captures the movement from one location to another.
\end{enumerate}

\section*{True / False}
\begin{enumerate}[label=\arabic*.]
\setcounter{enumi}{5}
\item \textbf{Answer:} False
\\ \textit{Options:} A) True; B) False
\\ \textit{Note:} Pongal is primarily a harvest festival celebrated in South India, centered around agricultural gratitude and not specifically focused on the exchange of gifts or sibling bonds.
\item \textbf{Answer:} False
\\ \textit{Options:} A) True; B) False
\\ \textit{Note:} The term 'Mahavira' actually means 'Great Hero' in Jainism, and while he is revered as an enlightened being, the phrase 'The Enlightened One' more accurately refers to the Buddha in Buddhism.
\item \textbf{Answer:} True
\\ \textit{Options:} A) True; B) False
\\ \textit{Note:} The statement is true because in a conditional proposition, the antecedent is the part that comes before the "if," and in this case, it specifies the condition that the Bees do not win their first game.
\item \textbf{Answer:} True
\\ \textit{Options:} A) True; B) False
\\ \textit{Note:} In Jaina traditions, caityavasis are indeed recognized as monks who reside in forested areas and dedicate themselves to ascetic practices, reflecting their commitment to spiritual purification and discipline.
\item \textbf{Answer:} False
\\ \textit{Options:} A) True; B) False
\\ \textit{Note:} The statement is false because SCNT stands for somatic cell nuclear transfer, not somatic cellular nuclear transmission, and it is primarily associated with cloning rather than genetic engineering.
\end{enumerate}

\section*{Long Form Response}
\begin{enumerate}[label=\arabic*.]
\setcounter{enumi}{10}
\item \textbf{Answer:} After the Bar Kochba revolt, Babylonia and Europe emerged as the two principal centers for Jewish development, each contributing significantly to Jewish philosophy and culture. Babylonia became a hub for the development of the Talmud, which synthesized Jewish law and ethics, influencing Jewish thought for centuries. In contrast, Europe, particularly during the Middle Ages, fostered a rich intellectual environment where Jewish philosophers like Maimonides engaged with and integrated Greek philosophy into Jewish theology, promoting rationalism and ethical discourse. The cultural impact of these centers is profound, as they shaped Jewish identity and resilience in the face of adversity, ensuring the continuity of Jewish thought and practice in diverse contexts. Together, Babylonia and Europe represent the dynamic evolution of Jewish intellectual tradition, bridging ancient texts with contemporary philosophical inquiries.
\\ \textit{Note:} correct because it accurately identifies Babylonia and Europe as the two central locations for Jewish intellectual and cultural development post-Bar Kochba, emphasizing their unique contributions to Jewish philosophy through the Talmud in Babylonia and the integration of Greek thought in European Jewish scholarship.
\item \textbf{Answer:} The ten-day New Year festival in Babylonian culture, known as Akitu, holds significant philosophical implications as it embodies the cyclical nature of life and the renewal of order. This festival celebrates the reinstatement of the divine king, reflecting the values of legitimacy, community, and the relationship between the human and the divine. The rituals performed during Akitu symbolize the triumph of order over chaos, aligning with the Babylonians' belief in a cosmos governed by divine law. Additionally, the emphasis on communal participation during this festival underscores the collective identity and social cohesion that were vital to Babylonian society. Overall, the Akitu festival serves as a profound expression of the Babylonians' worldview, illustrating their reverence for the cyclical renewal of life and the importance of maintaining harmony within their society.
\\ \textit{Note:} correct because it effectively connects the ten-day New Year festival's rituals and significance to the Babylonians' philosophical beliefs about order, legitimacy, and the divine, illustrating how these elements reflect the core values of their society.
\end{enumerate}

\vfill
\noindent\textit{End of Answer Key}
\end{document}